\section{Conclusion}
\label{sec:conclusion}

This study introduced \textit{MicroPhys-BMS}, an end-to-end PIML-KD framework for real-time SOC estimation of LFP batteries on automotive microcontrollers. The core conclusions are:

\begin{enumerate}
	\item \textbf{Flat OCV Plateau Observability:} The 12-D physical feature space $X_{\text{PIML}}$ ($\Delta V_{10}$, $\tau_{\text{log}}$, $P_{\text{ohm}}$, $\Phi_{\text{Arr}}$ with $E_a = 35.0\ \text{kJ/mol}$) resolves flat-plateau unobservability. The student Micro-MLP achieved MAE $= 0.454\%$ and RMSE $= 1.236\%$ ($R^2 = 0.9986$) at $25^\circ\text{C}$, outperforming its 300-tree LightGBM teacher (MAE $= 0.488\%$).
	\item \textbf{Cross-Dataset Generalization:} On Stanford LFP, the student recorded MAEs of $2.48\%$ (YX07) and $3.35\%$ (YX08), surpassing the Stanford baseline ($4.30\%$) \cite{che2025enhanced} by $42.3\%$ and $22.1\%$. On the independent CALCE A123 LFP cell \cite{calce2014a123} (different manufacturer, 18650 form), the student achieved MAE $= 0.173\%$ with $98.89\%$ CDF compliance within $\le 2.0\%$.
	\item \textbf{Comparative Evaluation:} Against PINN \cite{karniadakis2021physics} ($13{,}185$ params), Tabular Transformer \cite{chen2024transformer} ($13{,}899$ params), and 5-model LightGBM ensemble \cite{ke2017lightgbm} ($\approx 157{,}500$ params), the 961-parameter Micro-MLP achieved competitive accuracy (Stanford MAE $0.536\%$ vs. $0.394\%$--$1.188\%$; CALCE MAE $0.173\%$ vs. $0.132\%$--$0.231\%$) with $13.7\times$--$164\times$ fewer parameters. It is the only model satisfying automotive embedded constraints (Flash $< 16$~kB, ASIL-D).
	\item \textbf{ASIL-D Deployment:} The student compiled into \texttt{bms\_soc\_piml\_mlp3.h} requires $3.75\ \text{kB}$ Flash, $192\ \text{Bytes}$ SRAM, zero \texttt{malloc}, and $12.0\ \mu\text{s}$ per inference on ARM Cortex-M4 at $160\ \text{MHz}$, achieving $154.2\times$ latency acceleration over reference RF ensembles \cite{yahyaei2023embedded}.
\end{enumerate}

The complete source code, trained weights, and reproducible validation pipeline are available at \url{https://github.com/ha22da/PIML}. Future work will extend to online joint SOH/SOP co-estimation and integrate the deterministic C engine into commercial VCU energy management systems.
